% 拟定论文题目：智能体系统失败诊断与可靠恢复方法研究
% 本文件可直接替换 whu-thesis-latex-template/chapters/chapter1.tex。
% 参考文献使用 bibs/ref1.bib，正文引用格式遵循模板的 gbt7714-numerical 样式。

\chapter{绪论}\label{chap1}

\section{研究背景与意义}

近年来，大语言模型驱动的智能体逐渐由单轮文本生成扩展为能够感知环境、调用工具、规划步骤并持续执行的任务系统。此类系统已被用于软件开发、知识推理、图形界面操作、协同决策和云系统运维等场景。与传统的一次性模型推理不同，智能体任务通常由多个相互关联的步骤构成：早期步骤产生的信息或环境变化会成为后续步骤的输入，一处错误可能沿执行链条传播，并最终表现为任务失败。因此，系统最终失败并不等于最后一步出错，也不能仅凭最终答案反推出真正的责任位置。

智能体系统的可靠性问题可以概括为两个连续而不同的阶段。第一阶段是\textbf{失败诊断}：从完整执行轨迹中找出最可能导致失败的责任步骤，并给出可比较的责任排序。第二阶段是\textbf{任务恢复}：根据诊断结果选择合适的干预位置和范围，重新执行并验证任务是否恢复，同时避免破坏原本正确的状态。前者回答“哪里出了问题”，后者回答“如何把系统带回可接受状态”。诊断本身不会修改错误，恢复也不能脱离诊断盲目进行。

现有研究已经分别关注失败归因、错误检测、回溯推理和故障缓解，但两个阶段之间仍缺少稳定的输出衔接。一方面，细粒度步骤定位明显难于责任智能体识别，长轨迹、冗余信息和跨步骤关联都会干扰判断；另一方面，即使找到可疑步骤，单点删除、从头重试或固定范围回退也未必能够撤销由该步骤传播出的后续影响。研究能够驱动实际恢复的细粒度失败诊断方法，以及由诊断结果引导、带有验证保护的可靠恢复方法，对提高智能体系统的可解释性、可维护性和任务成功率具有重要意义。

基于上述认识，本文拟定题目为：

\begin{center}
  \textbf{智能体系统失败诊断与可靠恢复方法研究}
\end{center}

该题目刻画的是智能体系统中普遍存在的“执行失败后如何定位并恢复”问题，不把某一种工具、数据集或实验约束写入题目。具体的执行环境和受控条件作为方法实现与实验边界在后文说明，而不替代研究问题本身。

\section{相关研究现状}

\subsection{智能体任务失败归因与细粒度定位}

智能体失败诊断研究首先需要回答：一次多步骤任务失败后，哪一个参与者、哪一个执行步骤应当承担主要责任。2025年，Zhang等首次系统提出面向大语言模型多智能体系统的自动失败归因问题，并构建Who\&When数据集，覆盖127个多智能体系统及其细粒度失败标注。研究发现，最佳方法识别责任智能体的准确率为53.5\%，但定位决定性错误步骤的准确率仅为14.2\%~\cite{zhang2025whowhen}。这一差距说明，从系统级失败下钻到步骤级责任位置并非简单的分类细化，而是一个独立且困难的问题。

2026年，Sun等针对长日志和专业知识不足造成的定位困难，提出“先划定失败范围、再定位具体步骤”的两阶段框架。该方法先缩小候选范围，再进行步骤判断，在Who\&When数据集上的步骤级准确率相对既有方法最高提升24.27\%~\cite{sun2026scope}。这项工作表明，控制诊断范围能够降低长轨迹中的推理负担，但候选范围内的步骤仍需要可靠排序。

同年，Geng等面向多智能体代码生成提出FLKR，并构建专家标注的COFL基准。该方法通过行为编码、知识策略对齐和一致性评分，降低“答案路径不同但语义合理”对定位的干扰；相对提示式基线，其故障定位准确率最高提高14个百分点，Top-1准确率最高提高45个百分点~\cite{geng2026flkr}。该结果说明，仅依赖表面输出或固定解题路径容易误判，定位方法需要识别步骤的行为作用及其上下文。

随后，Chen等提出TraceElephant基准，强调开发者真实诊断时能够观察到的不仅是智能体输出，还包括输入、上下文和完整执行轨迹。实验显示，完整轨迹相对部分可观测轨迹可使归因准确率最高提高76.5\%~\cite{chen2026traceelephant}。这进一步证明，可观测证据的完整性直接影响失败归因质量，也说明完整轨迹中的有效证据提取是方法设计的关键。

2026年，Zhu等提出StepFinder，将执行日志编码为时序语义序列，再通过时序建模和注意力机制捕获轨迹演化及跨步骤依赖。该方法在Who\&When上优于基于大模型直接推理的步骤归因方法，并将相对最快大模型方法的推理时间降低79\%~\cite{zhu2026stepfinder}。这表明，显式建模时序演化和跨步骤关联，不仅可能提高定位效果，也有助于降低反复长上下文推理的成本。

上述研究形成了由“提出任务与数据集”到“缩小失败范围”“增强行为知识”“补全轨迹观测”再到“建模时序依赖”的清晰发展路径。现有结果共同证明了细粒度责任定位的研究价值，但仍存在三方面不足：第一，步骤级定位性能仍明显落后于粗粒度责任识别；第二，局部异常证据与跨步骤交互信息通常被分开处理，难以形成覆盖全部步骤的统一责任排序；第三，多数研究把定位准确率视为终点，尚未充分验证诊断排序能否直接改善后续任务恢复。

\subsection{智能体任务恢复与错误干预}

任务恢复研究关注智能体在中间错误、环境变化或协作失配发生后，能否回到有效状态并继续完成任务。2025年，Guo等提出SyncMind和SyncBench，用21个真实GitHub仓库中的24332个实例评估协同软件工程智能体从环境认知失配中恢复的能力。不同智能体的恢复表现从不高于3.33\%到不低于28.18\%不等，协作意愿不高于4.86\%~\cite{guo2025syncmind}。该研究揭示，即使模型具备较强的生成能力，面对已经变化的共享环境时仍缺乏稳定的状态感知和适应能力。

同年，Huang等研究含故障成员的多智能体协作韧性。实验显示，不同协作结构在故障条件下的性能下降差异显著，其中层次结构下降5.5\%，另外两种结构分别下降10.5\%和23.7\%；进一步引入Challenger和Inspector后，最多可恢复故障智能体造成的96.4\%错误~\cite{huang2025resilience}。这说明，系统结构、复核机制与干预方式会直接影响故障后的恢复能力。

在可逆推理方向，Zhao等提出ReAgent，使多个智能体在发现冲突时能够回到较早的有效推理状态，并采用局部和全局两种回滚协议。该方法在三个多跳问答基准上相对仅向前推理的基线取得约6\%的稳定提升~\cite{zhao2025reagent}。这项工作证明，错误发生后继续向前并非唯一选择，保留有效状态并进行回溯能够抑制早期错误传播。

同年，Wu等面向图形界面智能体提出BacktrackAgent，通过Verifier、Judger和Reflector检测错误、判断回退时机并反思后续动作。实验在Mobile3M和Auto-UI上同时提高了任务成功率与步骤准确率~\cite{wu2025backtrackagent}。该研究把错误检测和恢复执行连接起来，但其公开摘要未给出可跨基准统一比较的单一提升比例，因此本文只引用其定性结论，不虚构合并数字。

在云系统可靠性场景中，Chen等提出STRATUS，将检测、诊断和故障缓解组织为状态机，并形式化“事务式无回归”要求，只有恢复目标且不引入额外回归的方案才可被接受。该系统在AIOpsLab和ITBench上的故障缓解成功率至少达到现有SRE智能体的1.5倍~\cite{chen2025stratus}。这说明恢复方法不仅要“做出修改”，还必须验证修改是否真正解决目标问题并控制副作用。

相关研究已经证明回溯、复核、状态机和无回归验证能够提高智能体系统的韧性，但仍存在以下空缺：第一，许多方法依赖固定回退策略或特定场景的检测器，缺少由细粒度责任排序直接驱动的恢复候选选择；第二，现有工作较少区分“扩大干预范围”与“选择真正相关的依赖步骤”各自带来的收益；第三，目标恢复与非目标损害有时被分别报告，尚缺少贯穿候选生成、执行和验收的一致事务式边界；第四，失败归因与恢复研究大多各自优化，诊断输出对恢复成功率的实际价值仍需通过同候选池下的排序、打乱和Oracle对照来验证。

\subsection{研究现状小结}

综合上述两条研究线，可以得到三个判断。其一，失败归因与任务恢复并非人为构造的小众问题，而是在多智能体协作、软件工程、知识推理、图形界面操作和云系统运维等不同应用中反复出现的共同可靠性问题。其二，研究重点正由最终结果判断逐步转向完整轨迹、步骤级责任、时序依赖和可逆执行。其三，当前最明显的连接缺口是：诊断方法通常停留在“找错”，恢复方法通常从自身的错误检测或固定规则出发，二者缺少统一、可验证的输出桥梁。

因此，本文不再把研究目标限定为某一种工具调用，也不把“Top-1是否命中”作为全部价值，而是研究从\textbf{细粒度责任排序}到\textbf{诊断引导可靠恢复}的完整问题链条。前一问题提供可操作的责任先验，后一问题通过实际执行验证该先验是否具有下游价值。

\section{研究问题与总体思路}

围绕智能体系统任务失败后的诊断和恢复，本文提出两个相互衔接、又可独立评价的研究问题。

\textbf{RQ1：如何利用执行轨迹中的局部行为证据与跨步骤关联，对智能体任务失败进行细粒度责任定位？}

RQ1的输入是一次已失败任务的完整执行轨迹，输出是全部执行步骤的责任分数和全排序。输出全排序而不是只输出Top-1，是因为后续恢复可能受到执行代价、可逆性和风险约束，需要依次考察多个候选位置。该问题拟从局部行为模式和跨步骤交互关系两类互补证据出发，研究稳定的责任排序方法。

\textbf{RQ2：如何利用失败诊断结果生成、选择并验证恢复方案，提高智能体系统的任务恢复能力并降低错误干预风险？}

RQ2以RQ1的责任排序为诊断先验，但不把诊断正确直接等同于恢复成功。该问题关注三个连续动作：优先尝试哪个可疑位置、需要共同处理多大的相关范围、如何通过重新执行和无回归检查决定提交或撤销。RQ2的主结果是任务是否真实恢复，而不是另一个定位指标。

两个问题的关系可以概括为：

\begin{center}
  失败轨迹 $\longrightarrow$ 责任步骤全排序 $\longrightarrow$ 恢复候选与范围 $\longrightarrow$ 执行验证与提交/撤销。
\end{center}

RQ1与RQ2的收益分别计算，任何跨问题的累计提升都不能记为RQ2自身提升。RQ2必须相对同一恢复任务上的最强公平恢复基线独立报告。

\section{研究内容与方法设计}

\subsection{基于局部证据与交互依赖的责任排序}

针对RQ1，本文拟研究两个完整模块。

（1）\textbf{局部责任候选排序模块。} 对轨迹中每个步骤的匿名行为结构和局部上下文进行表示，通过成对比较证据估计各步骤相对承担责任的可能性，得到覆盖全部步骤的初始排序。该模块聚焦单步附近“发生了什么”，不依赖任务答案身份泄漏。

（2）\textbf{跨步骤交互依赖排序模块。} 将步骤之间的时序邻接、同角色连续行为、相同动作形态和相同结构模式表示为多通道交互关系，在轨迹图上聚合邻域信息并生成新的全排序。该模块聚焦“前一步如何影响后一步”。这些关系是可计算的交互代理，不直接声称为真实因果边。

两个模块通过仅使用训练侧内部验证数据学习的冻结融合规则组合。局部证据有助于识别直接异常，交互依赖有助于解释错误传播；组合的目标是使责任步骤在完整排序中稳定前移，而不是增加多个零散特征后只报告单一Top-1结果。

\subsection{诊断引导的依赖感知可靠恢复}

针对RQ2，本文同样只保留两个完整模块，控制方法复杂度和验证成本。

（1）\textbf{诊断引导的恢复候选排序模块。} 将RQ1输出的责任排序转化为恢复候选的尝试顺序，并结合训练侧可见的候选后验冻结排序规则。该模块回答“先恢复哪里”，其价值通过保持恢复候选池不变、只替换原生排序为固定顺序或随机打乱顺序进行检验。

（2）\textbf{依赖感知的事务式恢复模块。} 以候选责任步骤为锚点，生成由单步到相关后续步骤的多粒度恢复计划；执行前保留检查点，执行恢复计划后同时验证目标任务和非目标行为，满足条件则提交，否则撤销。依赖范围选择与事务式提交/撤销共同构成一个恢复模块：前者决定“改多大范围”，后者保证“错误方案不被永久接受”。它们不拆成多个独立创新点。

需要特别说明，“诊断引导—依赖感知—事务验证”可以作为便于理解的三个环节，但在正式方法计数中，RQ2只有上述两个模块：诊断排序模块，以及把依赖范围与事务验收耦合起来的恢复模块。

\section{实验设计与消融归因原则}

\subsection{RQ1评价设计}

RQ1采用覆盖全部步骤的排序指标，并以任务组为配对单位进行统计检验。主要比较包括：最强公平排序基线、仅局部模块、仅交互依赖模块、完整组合、删除任一模块、打乱依赖关系、随机依赖关系以及Oracle上界。模型训练、融合权重选择和阈值确定均在外层测试划分之外完成，最终使用task-group paired bootstrap报告置信区间。

消融表述遵循“一个对照回答一个问题”：仅局部模块检验单步行为证据是否有效；仅依赖模块检验跨步关系是否有效；完整组合相对两个单模块的增益检验互补性；删除和打乱对照检验收益是否依赖正确的模块内容，而非参数量或额外计算。

\subsection{RQ2评价设计}

RQ2以任务级恢复成功率为主指标，并同时记录非目标损害。正式比较至少包括：不恢复、最强公平恢复基线、仅诊断排序的单点恢复、等规模但不使用依赖身份的范围扩展、固定排序的依赖恢复、打乱诊断排序的依赖恢复、完整方法和Oracle。

RQ2的消融归因必须拆成以下三层：

（1）\textbf{整体有效性。} 完整方法相对最强公平恢复基线的差值，回答整个方案是否值得使用。

（2）\textbf{恢复范围贡献。} 在相同诊断排序下，将单点恢复与等规模范围扩展比较，回答“从一个动作扩展到多个动作”本身贡献多少。

（3）\textbf{依赖身份贡献。} 在尽量匹配恢复范围和预算后，将依赖感知范围与非依赖范围比较，回答“选择这些相关动作”是否比“仅仅多处理几个动作”更好。

因此，不允许把完整方法的全部提升归因于依赖语义。若范围扩展已经带来主要收益，只能声称依赖结构在其上提供额外增益；只有该边际比较的配对置信区间下界大于零时，才可进一步声称依赖身份具有独立且稳健的贡献。事务式验收在当前可行性阶段作为所有恢复条件共享的安全边界；若要声称其独立增益，正式论文还需在相同计划和排序下增加“仅验证目标”与“目标加无回归验证”的单一受控对照。

\section{前期可行性实验与初步结果}

前期实验已分别验证两个研究问题的指标可行性，但不将其表述为毕业论文的最终外部效度结论。

在RQ1上，局部责任候选排序相对最强公平基线取得9.8031\%的稳定正提升；跨步骤交互依赖模块取得33.3350\%的平均相对提升；完整组合取得44.3302\%的平均相对提升。完整组合的任务组配对bootstrap绝对差95\%置信区间为$[0.13237, 0.17164]$，下界大于零。删除、打乱和随机依赖对照均低于完整方法，说明局部证据与跨步骤关系具有互补价值。

在RQ2的34个受控可执行任务上，最强公平恢复基线恢复$7/34$个任务，完整的诊断引导依赖恢复方法恢复$14/34$个任务，相对提升100\%；任务级配对bootstrap绝对差95\%置信区间为$[0.08824, 0.35294]$。同时，等规模非依赖范围恢复$12/34$个任务，完整依赖范围在其上净增2个任务。故准确结论是：\textbf{诊断排序、前向范围扩展与相关步骤选择组成的完整方案有效；其中范围扩展贡献较大，依赖身份提供额外正信号，但当前样本尚不足以把全部收益归给依赖语义。}

上述结果达到当前开题阶段冻结的可行性门：两个研究问题的完整组合均相对各自最强公平基线取得不低于10\%的提升，且主比较的配对置信区间下界大于零。后续正式论文仍需扩大自然失败任务外测、补充强基线复现，并在不增加模块数量的前提下验证结论的跨场景稳定性。

\section{预期创新点}

本文的预期创新点概括如下：

（1）提出一种联合局部行为证据与跨步骤交互关系的智能体任务失败责任排序方法，从仅判断最终失败或责任智能体进一步下钻到覆盖全部步骤的细粒度责任排序。

（2）提出一种由失败诊断结果直接驱动的可靠恢复方法，将责任排序用于恢复候选选择，并通过相关范围恢复与事务式验证把“找错”连接到“实际恢复”。

（3）构建可证伪的诊断—恢复评价链路，在同一候选池下设置固定、打乱和Oracle排序对照，并将范围扩展贡献与依赖身份贡献分开报告，避免把组合收益错误归因于单一机制。

\section{论文组织结构}

全文拟分为六章。第一章介绍研究背景、相关研究现状、研究问题、总体方法与前期可行性证据。第二章介绍智能体执行轨迹、失败归因、排序学习、依赖分析与事务恢复的相关概念和技术基础。第三章研究基于局部证据与交互依赖的细粒度责任排序方法。第四章研究诊断引导的依赖感知可靠恢复方法。第五章开展跨数据、跨任务和跨场景实验，完成基线、消融、统计显著性、效率与威胁分析。第六章总结全文并讨论后续研究方向。
